\subsection{A flat reduced augmented ring can have a nonflat fixed locus}
\label{sec:kt-fixed-locus-nonflat}

The following elementary example addresses a specific missing
hypothesis in KT-NORMAL-FLAT-TRANSFER.  It is not asserted to
satisfy all the archived model conditions \texttt{M1--M4}.

\begin{lemma}[KT-FIXED-LOCUS-NONFLAT]
Let $O$ be a complete discrete valuation ring in which two
is invertible, with uniformizer $\varpi$.  For every $n\geq1$
there is a finite free, reduced, augmented local $O$-algebra
$R$ of rank three, with reduced generic fibre and an
augmentation-preserving involution $\kappa$, such that its
fixed-locus quotient $R/(r-\kappa(r):r\in R)$ has nonzero
$O$-torsion.  The odd cotangent at the augmentation has
length $n$.
\end{lemma}
\begin{proof}
Put $a=\varpi^n$ and, inside the product algebra $O^3$, set
\[
  1=(1,1,1),\qquad u=(0,a,-a),\qquad v=(0,a,a),
  \qquad R=O1\oplus Ou\oplus Ov.
\]
The determinant of the three displayed column vectors is
$2a^2\ne0$, so they are $O$-linearly independent.  Their
products satisfy
\[
                 u^2=av,\qquad uv=au,\qquad v^2=av.
\]
Thus $R$ is a subalgebra, finite free of rank three and
reduced.  Its generic fibre is the full product
$(\operatorname{Frac}O)^3$, since $2a^2$ becomes a unit.
Modulo $\varpi$, the span of $u,v$ is square-zero, so
$R/\varpi R$ is local with residue field $k(O)$.
Finite integrality over the local ring $O$ shows that $R$
itself is local; finite freeness makes it complete.

The first projection is an augmentation $x:R\to O$.
Interchanging the last two coordinates preserves $R$
and gives the involution
$\kappa(u)=-u$, $\kappa(v)=v$.
Because two is a unit, its fixed-locus ideal is $(u)$.
As an $O$-submodule this ideal is $Ou\oplus aOv$.
Consequently
\[
 R/(r-\kappa(r))\simeq O[v]/(av,v^2)
             \simeq O\oplus(O/a)v.
\]
The last summand is nonzero torsion, so the fixed locus
is not flat over $O$.  Finally $J=\ker x=Ou\oplus Ov$
and $J^2=aOu\oplus aOv$.  Hence
\[
                     (J/J^2)^-\simeq(O/a)u,
\]
whose length is $n$.
\end{proof}

The example has a base augmentation, a finite flat reduced
ambient algebra, and a reduced generic fibre.  These
properties alone therefore do not establish flatness of
the integral fixed-locus quotient or erase its odd normal
cotangent.  A source-specific saturation or integral
comparison theorem remains necessary.
